\documentclass[11pt,a4paper]{article}
\input{config}

\begin{document}
\titlepageandcontents

%---------------------------------------------------------------------------
\section{Zadání}

\begin{itemize}
\item Vytvořte implementaci dopředné a zpětné vlnkové transformace CDF 9/7 a CDF 5/3 pomocí lifting schématu.
\item CDF 5/3 implementujte bezeztrátově a v integerové aritmetice.
\item Vytvořte vizualizaci řešení -- aplikaci, která dokáže dopředně a zpětně transformovat obrázek ve formátu ppm.
\item Změřte výkonnost řešení a porovnejte s CPU implementací.
\item Popište způsob využití GPU.
\end{itemize}

%---------------------------------------------------------------------------
\section{Použité technologie}

\begin{itemize}
\item OpenCL 2.2
\end{itemize}

%---------------------------------------------------------------------------
\section{Použité zdroje}

\begin{itemize}
\item D. Bařina. \textit{Lifting Scheme Cores for Wavelet Transform.} PhD thesis, Brno University of Technology, Brno, 2015.
\end{itemize}

%---------------------------------------------------------------------------
\section{Způsob akcelerace}

Dvojrozměrná vlnková transformace CDF 5/3 byla implementována pomocí paralelního jádra z [Bařina 2015], tedy pomocí dvou kernelů.
Transformace CDF 9/7 byla implementována pomocí dvou sekvenčně aplikovaných paralelních jader stejné struktury, tedy pomocí čtyř kernelů.
Vícestupňový rozklad je proveden voláním stejného algoritmu na nízkofrekvenční složku transformace.
Pro akceleraci přístupu ke dvourozměrným datům je využita texturovací paměť.

%---------------------------------------------------------------------------
\section{Ovládání vytvořeného programu}

Program lze spustit pomocí:
\begin{verbatim}
./dwt {forward|inverse} {cpu|gpu} {53|97} <input.ppm> <output.ppm>
\end{verbatim}
Kde
\begin{itemize}
\item \texttt{forward}, resp. \texttt{inverse} značí dopřednou, resp. inverzní transformaci.
\item \texttt{cpu}, resp. \texttt{gpu} značí zavolání implementace na cpu, resp. na gpu.
\item \texttt{53}, resp. \texttt{97} značí provedení transformace CDF 5/3, resp. CDF 9/7.
\item \texttt{<input.ppm>}, resp. \texttt{<output.ppm>} je vstupní, resp. výstupní soubor ve formátu PPM.
\end{itemize}

%---------------------------------------------------------------------------
\section{Vyhodnocení}

Projekt byl vyhodnocen na stroji s následujícími parametry:
\begin{verbatim}
Intel(R) Core(TM) i5-8300H CPU @ 2.30GHz
GP107M [GeForce GTX 1050 Mobile]
DDR4 8192 MB SODIMM @ 2.4 GHz
\end{verbatim}

Měření bylo provedeno na náhodně vygenerovaných obrázcích (šum) čtvercových roměrů o straně $2^7$ pixelů až $2^{13}$ pixelů.
Pro každou kategorii byl z naměřených dat zvolen medián jako reprezentativní hodnota.
Na CPU byl měřen pouze samotný běh algoritmu, na GPU je do času zahrnut i přenos dat na jednotku a zpět.
Naměřená data jsou vizualizována v následujících grafech:

\includegraphics[width=.78\textwidth]{images/Forward53.pdf}

\includegraphics[width=.78\textwidth]{images/Forward97.pdf}

\includegraphics[width=.5\textwidth]{images/Inverse53.pdf}%
\includegraphics[width=.5\textwidth]{images/Inverse97.pdf}

Z grafů je zřejmé, že efektivita implementované akcelerace transformace roste s velikostí transformovaného obrázku.
To platí jak pro transformaci CDF 9/7, tak i 5/3.
Jelikož jde o symetrický algoritmus, inverzní transformace dle očekávání dosahuje stejných hodnot jako transformace dopředná.
Dle očekávání je transformace CDF 9/7 dvakrát pomalejší oproti CDF 5/3.
Je to způsobeno dvojnásobným množstvím operací pro transformaci jednoho pixelu.

%---------------------------------------------------------------------------
\section{Rozdělení práce v týmu}

Projekt byl řešen individuálně.

%---------------------------------------------------------------------------
\section{Doporučení pro budoucí zadávání projektů}

Projekt jsem nekonzultoval, takže nemůžu nic podotknout.

\end{document}
% vim:set ft=tex expandtab enc=utf8:
